\sect{7.1\quad 引言}
\src{1}
与有限拟阵 $M$ 相关联的若干单纯复形之间，有着十分优美的联系。它们以面数和 Betti 数承载 $M$ 的某些重要不变量，并产生有用的代数结构。本章研究三种复形：(1)独立子集组成的拟阵复形 $\IN(M)$；(2)相对于底集次序 $\omega$ 的断路复形 $\BC_\omega(M)$；(3)相应几何格 $L$ 中链所组成的序复形 $\Delta(\bar L)$。
为了系统研究这些复形的组合性质和同调性质，我们采用可壳性。如果一个复形的所有刻面等势，且能排成某种便于归纳论证的次序，就称它可壳。Provan (1977) 证明了拟阵复形和断路复形的可壳性，Björner (1980a) 证明了几何格序复形的可壳性。我们使用的一项关键性质是：可壳复形 $\Delta$ 对应一个非负整数系数多项式 $h_\Delta(x)$，它编码了 $\Delta$ 的基本组合不变量。$h_\Delta(1+\lambda)$ 的系数是面数，而 $h_\Delta(0)$ 是最高维 Betti 数，其余 Betti 数均为零。由于 $h_\Delta(x)$ 的每个系数都可解释为某些刻面的个数，一旦建立基本理论，可壳复形的同调计算便成为纯粹的组合任务。
与拟阵最自然地联系起来的单纯复形，是它的独立集族 $\IN(M)$。研究其可壳性，会自然引出基的内活性和外活性，进而引出二元生成函数——Tutte 多项式 $T_M(x,y)$。其中 $T_M(x,1)$ 和 $T_M(1,y)$ 分别是 $\IN(M)$ 与 $\IN(M^*)$ 的壳多项式，$M^*$ 是对偶拟阵。作为应用，我们将导出若干拟阵不等式。
拟阵 $M=M(S)$ 相对于 $S$ 的次序 $\omega$ 的断路复形 $\BC_\omega(M)$，由 $S$ 中不包含任何断路的子集组成；断路是从一个回路中删去其最小元素所得的集合。Whitney (1932)、Rota (1964)、Wilf (1976) 及 Brylawski (1977a) 最初为计数目的发展了这一概念。断路复形承载 $M$ 的“色多项式”性质：其壳多项式等于 $T_M(x,0)$，面数是第一类 Whitney 数，并且 $\BC_\omega(M)$ 是某个相关复形上的锥，该相关复形的最高维 Betti 数为 $M$ 的 beta 不变量 $\beta(M)$。
Folkman (1966) 的开创性工作确定了几何格复形 $\Delta(\bar L)$ 的同调，此后这一同调一直发挥重要作用。一方面，Folkman 消失定理使几何格成为 Cohen--Macaulay 偏序集理论的动机性例子之一（Baclawski 1980；Björner、Garsia 和 Stanley 1982）。另一方面，Orlik 和 Solomon (1980) 证明：复超平面配置的补空间的奇异上同调环，可完全由交格的序同调描述。因此，在这些联系以及 Gel'fand 和 Zelevinskii (1986) 等其他联系中，几何格同调与拟阵在数学中的有趣应用相关。
在 $\IN(M)$ 和 $\Delta(\bar L)$ 的同调研究中，断路复形起着有趣的作用：$\BC_\omega(M^*)$ 诱导的循环构成 $\IN(M)$ 同调的一个与域的特征无关的基，而 $\BC_\omega(M)$ 类似地确定 $\Delta(\bar L)$ 同调的一个基。这些例子连同 Orlik--Solomon 代数，表明断路思想在构造与拟阵、几何格相关的代数对象的基方面具有某种普遍性。
\src{2}
本章旨在对上述主题作统一、简洁而平易的介绍，尽量少要求预备知识。第7.2--7.6节完全是组合性的，代数内容均推迟到最后四节。第7.7节简要介绍所需的单纯同调，使本章基本自足。正文只发展最基本的思想，更多结果与分支见习题。第7.11节集中给出原始文献出处及有关评注。
